\documentclass[11pt,a4paper]{article}
\usepackage[T1]{fontenc}
\usepackage{isabelle,isabellesym}
\usepackage{amssymb}
\usepackage{underscore}
\usepackage{pdfsetup}

\urlstyle{rm}
\isabellestyle{it}

\begin{document}

\title{A Mechanised Core of the TRIEL Specification Language}
\author{Dmitri Chistyakov}
\maketitle

\begin{abstract}
This document is a core of the semantics of the TRIEL specification language,
formalised and checked in Isabelle/HOL. It is based only on the public description of
the language (\texttt{TECHNICAL\_REPORT.md}, Sections 2.5 and 2.9, and
\texttt{FOUNDATIONS.md}). All proofs are complete: there are no unproved lemmas and no
axioms beyond those of Isabelle/HOL.
\end{abstract}

\section*{What is formalised}

\paragraph{Expressions (theory \texttt{TRIEL\_Core}).}
States are partial maps from names to values, ordered by information (extension).
Boolean expressions are evaluated in strong Kleene three-valued logic; the value-presence
predicate \texttt{PRESENT(x)} is always defined. Proved: commutativity and associativity
of \texttt{AND} and \texttt{OR}, De Morgan's laws and double negation (T1); monotonicity of
every \texttt{PRESENT}-free expression under extension of the state, the claim of
Section~2.9 (T2); non-monotonicity of \texttt{PRESENT} (T3); agreement with two-valued
logic when every name read is defined (T4); the meaning of \texttt{PRESENT} (T5).

\paragraph{The indicator predicate (theory \texttt{TRIEL\_Ex}).}
Following the logics of quasiary predicates of Nikitchenko and Shkilniak, a predicate is
given by its truth and falsity domains, and the total indicator predicate $E_x$ is true
where $x$ has a value and false where it has none. Proved: \texttt{PRESENT(x)} is $E_x$
(\texttt{present\_is\_Ex}); $E_x$ is total and not equitone; $E_x$ holds exactly when the
value of $x$ equals some value, as for the existence predicate of free logic; every
\texttt{PRESENT}-free expression is equitone; the comparison $x = x$ is the partial
indicator predicate.

\paragraph{Trace semantics (theory \texttt{TRIEL\_Trace}).}
Traces are non-empty sequences of (state, timestamp, event) entries with non-decreasing
timestamps; executions end \texttt{Done}, \texttt{Interrupted} or \texttt{Violated}. The
denotations of \texttt{MUST}, \texttt{MAY}, \texttt{MUST\_NOT}, \texttt{THEN},
\texttt{AND}, \texttt{OR} and \texttt{UNLESS} follow Section~2.5. Proved: the laws of
\texttt{UNLESS} --- \texttt{t UNLESS false DO e = t}, distributivity over \texttt{OR},
distributivity over \texttt{THEN} for a total right operand, and idempotence; a
counterexample showing that distributivity over \texttt{THEN} fails without totality; the
inequality of nested \texttt{UNLESS} and \texttt{UNLESS} with a disjunctive guard.

\paragraph{Terms (theory \texttt{TRIEL\_Terms}).}
Terms form an inductive type with a compositional denotation. Proved: every term is total
(\texttt{den\_total}), and therefore all four laws of \texttt{UNLESS}, including
distributivity over \texttt{THEN}, hold for every term without side conditions.

\paragraph{Not formalised.}
Arithmetic, further comparisons, the type system, staleness, deadlines, breach handling,
the conditional and event-triggered terms, and the satisfaction of invariants.

\newpage
\tableofcontents
\newpage

\parindent 0pt\parskip 0.5ex

\input{session}

\end{document}
